\documentclass[xcolor=table]{beamer}
\usepackage{beamerthemesplit}
\usepackage{wrapfig}
\usetheme{SPbGU}
\usepackage{pdfpages}
\usepackage{amsmath}
\usepackage{cmap}
\usepackage[T2A]{fontenc}
\usepackage[utf8]{inputenc}
\usepackage[english]{babel}
\usepackage{indentfirst}
\usepackage{amsmath}
\usepackage{tikz}
\usepackage{multirow}
\usepackage[noend]{algpseudocode}
\usepackage{algorithm}
\usepackage{algorithmicx}
\usepackage{fancyvrb}
%usepackage{fancyvrb}
%\usepackage{minted}
%\usepackage{verbments}


\beamertemplatenavigationsymbolsempty

\title[САГ $\rightarrow$ СМУ]{От синтаксического анализа графов к системам матричных уравнений}
\institute[СПбГУ]{
JetBrains Research, Programming Languages and Tools Lab  \\
Санкт-Петербургский Государственный Университет
}

\author[Юлия Сусанина]{Юлия Сусанина}

\date{14.12.2019}

\begin{document}
{
\begin{frame}[fragile]
  \begin{tabular}{p{2.0cm} p{7.5cm} p{1cm}}
   \begin{center}
      \includegraphics[height=1.5cm]{pictures/jetbrainsResearch.pdf}
    \end{center}
    &
    \begin{center}
      \includegraphics[height=1.5cm]{pictures/YC_logo.pdf}
    \end{center}
    &
    \begin{center}
      \includegraphics[height=1.5cm]{pictures/SPbGU_Logo.png}
    \end{center}
  \end{tabular}
  \titlepage
\end{frame}
}

\begin{frame}[fragile]
  \transwipe[direction=90]
  \frametitle{Context-Free Path Querying (CFPQ)}
  
\begin{itemize}
% \item Постановка задачи:
%  \begin{itemize}
  \item КС-грамматика $G=(N, \Sigma, R)$ \\
  $\mathcal{L}(G_S) = \{\omega \mid S \Rightarrow_{G}^{*} \omega\}$, $S \in N$
  \pause
  \item Орграф $D = (V,E,\sigma)$, $\sigma \subseteq \Sigma$, $E\subseteq V\times \sigma \times V$ \\
  $m\lambda n$ --- путь из $m$ в $n$ в графе $D$, $\lambda$ --- слово данного пути
  \pause
  \item $R_A = \{(m, n) \mid m\lambda n \text{ --- путь в } D, \lambda \in \mathcal{L}(G_A)\}$
  \pause 
  
  \vspace{-9pt}
  \begin{tikzpicture}
  \draw[line width=0.25mm, red] (0,0) -- (7.3,0)
  \end{tikzpicture}
%  \end{itemize}
%   \pause
%   \item Области примерения:
%   \begin{itemize}
%       \item графовые базы данных
%       \item биоинформатика
%       \item статический анализ кода
%   \end{itemize}
\end{itemize}
\end{frame}

\begin{frame}[fragile]
  \transwipe[direction=90]
  \frametitle{Вычислительные методы}
\begin{itemize}
  \item Использование достижений линейной алгебры
  \item Приближенные методы для ускорения вычислений
  \item Развитие и постоянное улучшение библиотек и программных пакетов, связанное с развитием искусственного интеллекта 
  \item Возможность применения параллельных вычислений
\end{itemize}
\end{frame}

\begin{frame}[fragile]
  \transwipe[direction=90]
  \frametitle{Цель}
\begin{itemize}
  \item Создание подхода, основанного на методах линейной алгебры и вычислительной математики, к задаче CFPQ
\end{itemize}

% Задачи:

% \begin{itemize}
%   \item сведение CFPQ к задаче решения (систем) матричных уравнений
%   \item исследование способов решения и свойств данного подхода
%   \item реализация и экспериментальное исследование
% \end{itemize}

\end{frame}


\begin{frame}[fragile]
  \transwipe[direction=90]
  \frametitle{Сведение CFPQ к решению матричных уравнений (1)}
  
  $$S \rightarrow aSb \mid ab$$
  $$T_E \in \mathbb{M}^{|V| \times |V|}: (T_{E})_{ij} = 1 \iff (i,j) \in R_E \; \forall E \in (N \cup \Sigma)$$
  \pause
  $$\Downarrow$$
    \centering
    \begin{tabular}{cl}
    \parbox{0.07\linewidth}{
    $$\{ T_S^k \}:$$
    }
    & \begin{tabular}{l}
    $T_S^0 = \mathbf{0}$ \\
    $T_S^{k+1} = T_a T_S^k T_b + T_a T_b$
    \end{tabular} \\
    \end{tabular}

  $$T_S^{\infty} - \texttt{наименьшее решение } T_S = T_aT_ST_b + T_aT_b$$

\end{frame}

\begin{frame}[fragile]
  \transwipe[direction=90]
  \frametitle{Сведение CFPQ к решению матричных уравнений (2)}
  
    $$\Downarrow$$
    \centering
    \begin{tabular}{cl}
    \parbox{0.07\linewidth}{
    $\{ \mathcal{T}_S^k \}:$
    }
    & \begin{tabular}{l}
    $\mathcal{T}_S^0 = \mathbf{0}$ \\
    $\mathcal{T}_S^{k+1} = \epsilon(T_a \mathcal{T}_S^k T_b + T_a T_b)$
    \end{tabular} \\
    \end{tabular}

    \vspace{20pt}
    \(
    \left.
    \begin{array}{c}
    \mathcal{T}_S^{\infty} - \texttt{наименьшее решение } \mathcal{T}_S = \epsilon(T_a \mathcal{T}_S T_b + T_a T_b), \\
    \texttt{где } \epsilon \texttt{ т.ч. } \mathcal{T}_{S}^{k} \leq \textbf{1} \quad \forall k
    \end{array}
    \right.
    \)
  
  \pause
  
  $$(\mathcal{T}_S^{k+1})_{ij} > 0 \iff (T_S^{k+1})_{ij} = 1$$
  $$ceil(\mathcal{T}_S^{\infty}) = T_S^{\infty}$$


\end{frame}

\begin{frame}[fragile]
  \transwipe[direction=90]
  \frametitle{Способы решения}
 
\begin{itemize}
    \item Линейный случай
    \begin{itemize}
        \item линейные уравнения специального типа \\
        (уравнения Сильвестра: $AXB + CXD = F$)
        \item линейные системы $Ax=b$
    \end{itemize}
    \item Нелинейный случай
    \begin{itemize}
        \item метод Ньютона \\
                $$X = G(X) \Rightarrow F(X) = X - G(X) = \mathbf{0}$$
                \(
                \left.
                \begin{array}{l}
                X_{i+1} = X_i - (F'(X_i))^{-1}F(X_i)
                \end{array}
                \right.
                \iff
                \left\{
                \begin{array}{l}
                F'(X_i)H_i = - F(X_i) \\
                X_{i+1} = X_i + H_i
                \end{array}
                \right.
                \)
    \end{itemize}
\end{itemize}
\end{frame}

\begin{frame}[fragile]
  \transwipe[direction=90]
  \frametitle{Первая реализация}

\begin{itemize}
    \item \textit{SciPy}
    \begin{itemize}
        \item sSLV --- решение разреженной системы линейных уравнений
        \item dNWT --- нахождение корней уравнения, метод Ньютона
    \end{itemize}
    \item Результаты сравнения с матричным алгоритмом (в мс)

        \begin{table}[h]
        \centering
        \begin{tabular}{ | c | c || c | c || c | c | c |}
        \hline
        Ontology    & |V| & dNWT & sSLV & dGPU & sCPU &  sGPU \\
        \hline
        \hline
        bio-meas    & 341 &  284 & 35   & 276  & 91  & 24\\
        people-pets & 337 &  73  & 49   & 144  & 38  & 6\\
        funding     & 778 &  502 & 184  & 1246 & 344 & 27\\
        wine        & 733 &  791 & 171  & 722  & 179 & 6\\
        pizza       & 671 &  334 & 161  & 943  & 256 & 23\\
        \hline
        \end{tabular}
        
        \end{table}

\end{itemize}
\end{frame}

\begin{frame}[fragile]
  \transwipe[direction=90]
  \frametitle{Результаты}
  
\begin{itemize}

  \item \textbf{ACM SIGMOD 2020 Student Research Competition:} \\ 
  Yuliya Susanina. Context-Free Path Querying via Matrix Equations.\\
  \textit{(принят во второй тур соревнования)}
  
\vspace{10pt}\hspace{-30pt} (+) \textit{Результаты работы в первом полугодии}
  
  \item \textbf{CIBB 2019:} \\
  Yuliya Susanina, Anna Yaveyn and Semyon Grigorev. Modification of Valiant’s Parsing Algorithm for String-Searching Problem. \\
  \textit{(доклад, ожидается публикация)}
  \item \textbf{Журнал <<Труды ИСП РАН>>:} \\
  Сусанина Ю.А., Григорьев С.В. Модификация алгоритма Валианта для задачи поиска подстрок. \\
  \textit{(ожидается публикация)}
  
\end{itemize}

\end{frame}




\begin{frame}[fragile]
  \transwipe[direction=90]
  \frametitle{Дальнейшие планы}

\begin{itemize}
    \item Эффективная реализация предложенного подхода с использованием специализированных библиотек и параллельных вычислений
    \item Определение подклассов полиномиальных уравнений, решение которых может быть сведено к CFPQ
    \item Попытка построить взаимное сведение между CFPQ и решением соответсвующих подклассов уравнений
    \item Планируется публикация результатов на GRADES-NDA 2020
\end{itemize}

\end{frame}




\end{document}